\section{Introduction}
\subsection{Why?}

在过去两个多世纪里，天文学家们已经收集了太阳系中约2亿次小行星的天体测量观测数据 。这个天体测量数据集对于估算近地小行星等太阳系天体（SSOs）的运行轨道至关重要 。用于SSOs轨道确定的绝大多数天体测量数据都是光学的，即在特定历元下的赤经（RA）和赤纬（DEC）成对测量值 。相比之下，只有极小一部分的SSO曾被雷达观测过 。因此，提高光学天体测量的质量是一项非常有价值的工作，特别是当这意味着可以延长对感兴趣天体的观测弧长时 

在过去两个多世纪里，天文学家积累了约两亿条小行星天体测量观测数据。这一庞大的历史数据集是精确确定太阳系天体（SSO）运行轨道，尤其是近地小行星轨道的基石。然而，这些宝贵的历史数据在归算过程中使用了种类繁多的星表作为参考框架，导致了不可避免的系统性误差。这些源于不同星表的局部偏差会直接影响轨道拟合的精度，进而降低星历预报的可靠性。

\subsection{Who?}
为了解决这一问题，天文学界持续致力于识别并修正这些系统性误差。早期的关键研究，如Chesley et al. (2010)利用2MASS星表识别了USNO系列星表中的局部偏差，以及Farnocchia et al. (2015)通过PPM XL星表对几乎所有主流星表进行了更广泛的去偏差处理，均取得了显著成效。这些工作极大地改善了小行星的轨道拟合后残差。

\subsection{What?}
本报告的工作继承并发展了这一思路，其核心目标是利用欧洲空间局盖亚（Gaia）任务第二次数据发布（DR2）所提供的、精度前所未有的天体测量数据作为“黄金标准”，对小行星天体测量中广泛使用的26个星表进行系统性偏差的重新量化与修正。本报告将首先介绍盖亚DR2作为新一代天体测量基准的卓越性能，随后详细阐述我们的修正方法论，展示对关键星表的分析结果，并评估新修正方案对小行星星历预报精度的实际影响。

\subsection{How?}
本研究的最终贡献是，为学界提供了一套基于盖亚DR2黄金标准的、迄今最高质量的星表系统性偏差修正表，从而为彻底挖掘海量历史观测数据的全部潜力提供了决定性工具。
\newpage